\scanpage{359}\mspage{275}
On a alors ($f$ \unsure{étant|ét.} $\pi_0$-accessible)
\[
 f(x)=\underset I{\underrightarrow{\lim}}\,f(x_i),
\]
les $f(x_i)$ sont dans $N_{\pi_1}$, donc dans $N_\pi$,
or $N_\pi$ est stable par $\unsure{\underline{\lim}|\underrightarrow{\lim}}$ de
cardinal $\leq\pi$ (puisque $N$ est engendrée
par $N_{\pi_1}$, formé\unsure{e} d'objets $\pi_1$ accessibles
et a fortiori $\pi$-accessibles, cf
SGA 4 I 9.9. D'où $f(x)\in N_\pi$ \unsure{c.q.f.d.}

\uline{Corollaire} 4.12.2. \inserted{Supposons que $\pi_0$ soit utile pour $N$, et que $\pi_1\geq\operatorname{Card\,red} N_{\pi_0}$.} Soit $\pi$ un cardinal $\geq\pi_1$,
et tel que $\pi^{\pi_0}=\pi$. Alors $f$ \deleted{est isom}
s'identifie, à équivalences de catégories
et à isom \inserted{de foncteurs} près, au foncteur
\[
 \mathrm{Ind}(f_\pi):\mathrm{Ind}_\pi(M_\pi)\unsure{-|\longrightarrow}\mathrm{Ind}_\pi(N_\pi),
\]
où $f_\pi:M_\pi\longrightarrow N_\pi$ est le foncteur
induit par $f$.

% slanted
\margin{\unsure{remonte|remonter} à 4.7.}
\uline{Proposition} 4.12.3. Soient $C$ une catégorie
équivalente à une petite catégorie, $N$
une \inserted{\unsure{$\mathcal U$-|$\mathscr U$-}}catégorie, $\pi_0$ un cardinal infini.
% archive sheet number: 358
